\documentclass{amsart}
% For dealing with references we use the comment environment
\usepackage{verbatim}
\newenvironment{reference}{\comment}{\endcomment}
%\newenvironment{reference}{}{}
\newenvironment{slogan}{\comment}{\endcomment}
\newenvironment{history}{\comment}{\endcomment}

% For commutative diagrams we use Xy-pic
\usepackage[all]{xy}

% We use 2cell for 2-commutative diagrams.
\xyoption{2cell}
\UseAllTwocells

% We use multicol for the list of chapters between chapters
\usepackage{multicol}

% This is generally recommended for better output
\usepackage{lmodern}
\usepackage[T1]{fontenc}

% For cross-file-references

% Package for hypertext links:
\usepackage{hyperref}

% For any local file, say "hello.tex" you want to link to please

% Theorem environments.
%
\theoremstyle{plain}
\newtheorem{theorem}[subsection]{Theorem}
\newtheorem{proposition}[subsection]{Proposition}
\newtheorem{lemma}[subsection]{Lemma}

\theoremstyle{definition}
\newtheorem{definition}[subsection]{Definition}
\newtheorem{example}[subsection]{Example}
\newtheorem{exercise}[subsection]{Exercise}
\newtheorem{situation}[subsection]{Situation}

\theoremstyle{remark}
\newtheorem{remark}[subsection]{Remark}
\newtheorem{remarks}[subsection]{Remarks}

\numberwithin{equation}{subsection}

% Macros
%
\def\lim{\mathop{\mathrm{lim}}\nolimits}
\def\colim{\mathop{\mathrm{colim}}\nolimits}
\def\Spec{\mathop{\mathrm{Spec}}}
\def\Hom{\mathop{\mathrm{Hom}}\nolimits}
\def\Ext{\mathop{\mathrm{Ext}}\nolimits}
\def\SheafHom{\mathop{\mathcal{H}\!\mathit{om}}\nolimits}
\def\SheafExt{\mathop{\mathcal{E}\!\mathit{xt}}\nolimits}
\def\Sch{\mathit{Sch}}
\def\Mor{\mathop{\mathrm{Mor}}\nolimits}
\def\Ob{\mathop{\mathrm{Ob}}\nolimits}
\def\Sh{\mathop{\mathit{Sh}}\nolimits}
\def\NL{\mathop{N\!L}\nolimits}
\def\CH{\mathop{\mathrm{CH}}\nolimits}
\def\proetale{{pro\text{-}\acute{e}tale}}
\def\etale{{\acute{e}tale}}
\def\QCoh{\mathit{QCoh}}
\def\Ker{\mathop{\mathrm{Ker}}}
\def\Im{\mathop{\mathrm{Im}}}
\def\Coker{\mathop{\mathrm{Coker}}}
\def\Coim{\mathop{\mathrm{Coim}}}

% Boxtimes
%
\DeclareMathSymbol{\boxtimes}{\mathbin}{AMSa}{"02}

%
% Macros for moduli stacks/spaces
%
\def\QCohstack{\mathcal{QC}\!\mathit{oh}}
\def\Cohstack{\mathcal{C}\!\mathit{oh}}
\def\Spacesstack{\mathcal{S}\!\mathit{paces}}
\def\Quotfunctor{\mathrm{Quot}}
\def\Hilbfunctor{\mathrm{Hilb}}
\def\Curvesstack{\mathcal{C}\!\mathit{urves}}
\def\Polarizedstack{\mathcal{P}\!\mathit{olarized}}
\def\Complexesstack{\mathcal{C}\!\mathit{omplexes}}
% \Pic is the operator that assigns to X its picard group, usage \Pic(X)
% \Picardstack_{X/B} denotes the Picard stack of X over B
% \Picardfunctor_{X/B} denotes the Picard functor of X over B
\def\Pic{\mathop{\mathrm{Pic}}\nolimits}
\def\Picardstack{\mathcal{P}\!\mathit{ic}}
\def\Picardfunctor{\mathrm{Pic}}
\def\Deformationcategory{\mathcal{D}\!\mathit{ef}}

\setlength{\emergencystretch}{3em}
\begin{document}
\title[Stacks Project, Tag 07J6]{331 --- Every crystalline module crystal carries a canonical integrable connection}
\maketitle
\noindent Copyright (C) 2005--2025 Johan de Jong. Stacks Project authors. Source: \href{https://stacks.math.columbia.edu/tag/07J6}{Tag 07J6}.

\noindent Editorial difficulty note (not author text). Difficulty H2: The connection must be constructed functorially from the crystal isomorphisms on a first-order divided-power thickening. Integrability is a separate substantive obligation: the cocycle identity on a specifically constructed second-order thickening must be expanded and identified with the curvature in the exterior square..

\noindent Editorial provenance note (not author text). History: excerpt from revision \texttt{a0444\allowbreak 6e57e\allowbreak c1fbc\allowbreak 252a8\allowbreak 71afc\allowbreak ec775\allowbreak 2fb28\allowbreak 07b14}, crystalline.tex, lines 2597--2690. Reference presentation was adapted; original-author context and core support blocks were retained or added. The mathematical wording is unchanged. Licensed under GNU FDL 1.2 or later; no invariant sections or cover texts. See ../licenses/COPYING. Context blocks reproduce author definitions and supporting results; they are not additional counted problems.

\begin{remark}
\label{remark-first-order-thickening}
In Situation \ref{situation-global}.
Let $(U, T, \delta)$ be an object of $\text{Cris}(X/S)$.
Write $\Omega_{T/S, \delta} = (\Omega_{X/S})_T$, see
Lemma \href{https://stacks.math.columbia.edu/tag/07IZ}{lemma module of differentials (Tag 07IZ)}.
We explicitly describe a first order thickening $T'$ of
$T$. Namely, set
$$
\mathcal{O}_{T'} = \mathcal{O}_T \oplus \Omega_{T/S, \delta}
$$
with algebra structure such that $\Omega_{T/S, \delta}$ is an
ideal of square zero. Let $\mathcal{J} \subset \mathcal{O}_T$
be the ideal sheaf of the closed immersion $U \to T$. Set
$\mathcal{J}' = \mathcal{J} \oplus \Omega_{T/S, \delta}$.
Define a divided power structure on $\mathcal{J}'$ by setting
$$
\delta_n'(f, \omega) = (\delta_n(f), \delta_{n - 1}(f)\omega),
$$
see Lemma \href{https://stacks.math.columbia.edu/tag/07HH}{lemma divided power first order thickening (Tag 07HH)}.
There are two ring maps
$$
p_0, p_1 : \mathcal{O}_T \to \mathcal{O}_{T'}
$$
The first is given by $f \mapsto (f, 0)$ and the second by
$f \mapsto (f, \text{d}_{T/S, \delta}f)$. Note that both are compatible
with the divided power structures on $\mathcal{J}$ and $\mathcal{J}'$
and so is the quotient map $\mathcal{O}_{T'} \to \mathcal{O}_T$.
Thus we get an object $(U, T', \delta')$ of $\text{Cris}(X/S)$
and a commutative diagram
$$
\xymatrix{
& T \ar[ld]_{\text{id}} \ar[d]^i \ar[rd]^{\text{id}} \\
T & T' \ar[l]_{p_0} \ar[r]^{p_1} & T
}
$$
of $\text{Cris}(X/S)$ such that $i$ is a first order thickening whose ideal
sheaf is identified with $\Omega_{T/S, \delta}$ and such that
$p_1 - p_0 : \mathcal{O}_T \to \mathcal{O}_{T'}$
is identified with the universal derivation $\text{d}_{T/S, \delta}$
composed with the inclusion $\Omega_{T/S, \delta} \to \mathcal{O}_{T'}$.
\end{remark}

\begin{remark}
\label{remark-second-order-thickening}
In Situation \ref{situation-global}.
Let $(U, T, \delta)$ be an object of $\text{Cris}(X/S)$.
Write $\Omega_{T/S, \delta} = (\Omega_{X/S})_T$, see
Lemma \href{https://stacks.math.columbia.edu/tag/07IZ}{lemma module of differentials (Tag 07IZ)}.
We also write $\Omega^2_{T/S, \delta}$ for its second exterior
power. We explicitly describe a second order thickening $T''$ of $T$.
Namely, set
$$
\mathcal{O}_{T''} =
\mathcal{O}_T \oplus \Omega_{T/S, \delta} \oplus \Omega_{T/S, \delta}
\oplus \Omega^2_{T/S, \delta}
$$
with algebra structure defined in the following way
$$
(f, \omega_1, \omega_2, \eta) \cdot
(f', \omega_1', \omega_2', \eta') =
(ff', f\omega_1' + f'\omega_1, f\omega_2' + f'\omega_2,
f\eta' + f'\eta + \omega_1 \wedge \omega_2' + \omega_1' \wedge \omega_2).
$$
Let $\mathcal{J} \subset \mathcal{O}_T$
be the ideal sheaf of the closed immersion $U \to T$. Let
$\mathcal{J}''$ be the inverse image of $\mathcal{J}$ under the
projection $\mathcal{O}_{T''} \to \mathcal{O}_T$.
Define a divided power structure on $\mathcal{J}''$ by setting
$$
\delta_n''(f, \omega_1, \omega_2, \eta) =
(\delta_n(f), \delta_{n - 1}(f)\omega_1, \delta_{n - 1}(f)\omega_2,
\delta_{n - 1}(f)\eta + \delta_{n - 2}(f)\omega_1 \wedge \omega_2)
$$
see Lemma \href{https://stacks.math.columbia.edu/tag/07HI}{lemma divided power second order thickening (Tag 07HI)}.
There are three ring maps
$q_0, q_1, q_2 : \mathcal{O}_T \to \mathcal{O}_{T''}$
given by
\begin{align*}
q_0(f) & = (f, 0, 0, 0), \\
q_1(f) & = (f, \text{d}f, 0, 0), \\
q_2(f) & = (f, \text{d}f, \text{d}f, 0)
\end{align*}
where $\text{d} = \text{d}_{T/S, \delta}$.
Note that all three are compatible with the divided power structures
on $\mathcal{J}$ and $\mathcal{J}''$. There are three ring maps
$q_{01}, q_{12}, q_{02} : \mathcal{O}_{T'} \to \mathcal{O}_{T''}$
where $\mathcal{O}_{T'}$ is as in Remark \ref{remark-first-order-thickening}.
Namely, set
\begin{align*}
q_{01}(f, \omega) & = (f, \omega, 0, 0), \\
q_{12}(f, \omega) & =
(f, \text{d}f, \omega, \text{d}\omega), \\
q_{02}(f, \omega) & = (f, \omega, \omega, 0)
\end{align*}
These are also compatible with the given divided power
structures. Let's do the verifications for $q_{12}$: Note
that $q_{12}$ is a ring homomorphism as
\begin{align*}
q_{12}(f, \omega)q_{12}(g, \eta) & =
(f, \text{d}f, \omega, \text{d}\omega)(g, \text{d}g, \eta, \text{d}\eta) \\
& =
(fg, f\text{d}g + g \text{d}f, f\eta + g\omega,
f\text{d}\eta + g\text{d}\omega + \text{d}f \wedge \eta +
\text{d}g \wedge \omega) \\
& = q_{12}(fg, f\eta + g\omega) = q_{12}((f, \omega)(g, \eta))
\end{align*}
Note that $q_{12}$ is compatible with divided powers because
\begin{align*}
\delta_n''(q_{12}(f, \omega)) & =
\delta_n''((f, \text{d}f, \omega, \text{d}\omega)) \\
& =
(\delta_n(f), \delta_{n - 1}(f)\text{d}f, \delta_{n - 1}(f)\omega,
\delta_{n - 1}(f)\text{d}\omega + \delta_{n - 2}(f)\text{d}(f) \wedge \omega)
\\
& = q_{12}((\delta_n(f), \delta_{n - 1}(f)\omega)) =
q_{12}(\delta'_n(f, \omega))
\end{align*}
The verifications for $q_{01}$ and $q_{02}$ are easier.
Note that $q_0 = q_{01} \circ p_0$, $q_1 = q_{01} \circ p_1$,
$q_1 = q_{12} \circ p_0$, $q_2 = q_{12} \circ p_1$,
$q_0 = q_{02} \circ p_0$, and $q_2 = q_{02} \circ p_1$.
Thus $(U, T'', \delta'')$ is an object of $\text{Cris}(X/S)$
and we get morphisms
$$
\xymatrix{
T''
\ar@<2ex>[r]
\ar@<0ex>[r]
\ar@<-2ex>[r]
&
T'
\ar@<1ex>[r]
\ar@<-1ex>[r]
&
T
}
$$
of $\text{Cris}(X/S)$ satisfying the relations described above.
In applications we will use $q_i : T'' \to T$ and
$q_{ij} : T'' \to T'$ to denote the morphisms associated to the
ring maps described above.
\end{remark}

\begin{situation}
\label{situation-global}
Here $p$ is a prime number and $(S, \mathcal{I}, \gamma)$ is a divided power
scheme over $\mathbf{Z}_{(p)}$. We set $S_0 = V(\mathcal{I}) \subset S$.
Finally, $X \to S_0$ is a morphism of schemes such that $p$ is
locally nilpotent on $X$.
\end{situation}

\begin{lemma}
\label{lemma-automatic-connection}
In Situation \ref{situation-global}.
Let $\mathcal{F}$ be a crystal in $\mathcal{O}_{X/S}$-modules
on $\text{Cris}(X/S)$. Then $\mathcal{F}$ comes equipped with a
canonical integrable connection.
\end{lemma}

\begin{proof}
Say $(U, T, \delta)$ is an object of $\text{Cris}(X/S)$.
Let $(U, T', \delta')$ be the infinitesimal thickening of $T$
by $(\Omega_{X/S})_T = \Omega_{T/S, \delta}$
constructed in Remark \ref{remark-first-order-thickening}.
It comes with projections $p_0, p_1 : T' \to T$
and a diagonal $i : T \to T'$. By assumption we get
isomorphisms
$$
p_0^*\mathcal{F}_T \xrightarrow{c_0}
\mathcal{F}_{T'} \xleftarrow{c_1}
p_1^*\mathcal{F}_T
$$
of $\mathcal{O}_{T'}$-modules. Pulling $c = c_1^{-1} \circ c_0$
back to $T$ by $i$ we obtain the identity map
of $\mathcal{F}_T$. Hence if $s \in \Gamma(T, \mathcal{F}_T)$
then $\nabla(s) = p_1^*s - c(p_0^*s)$ is a section of
$p_1^*\mathcal{F}_T$ which vanishes on pulling back by $i$. Hence
$\nabla(s)$ is a section of
$$
\mathcal{F}_T
\otimes_{\mathcal{O}_T}
\Omega_{T/S, \delta}
$$
because this is the kernel of $p_1^*\mathcal{F}_T \to \mathcal{F}_T$
as $\mathcal{O}_{T'} = \mathcal{O}_T \oplus \Omega_{T/S, \delta}$
by construction. It is easily verified that $\nabla(fs) =
f\nabla(s) + s \otimes \text{d}(f)$ using the description of
$\text{d}$ in Remark \ref{remark-first-order-thickening}.

\medskip\noindent
The collection of maps
$$
\nabla : \Gamma(T, \mathcal{F}_T) \to
\Gamma(T, \mathcal{F}_T \otimes_{\mathcal{O}_T} \Omega_{T/S, \delta})
$$
so obtained is functorial in $T$ because the construction of $T'$
is functorial in $T$. Hence we obtain a connection.

\medskip\noindent
To show that the connection is integrable we consider the
object $(U, T'', \delta'')$ constructed in
Remark \ref{remark-second-order-thickening}.
Because $\mathcal{F}$ is a sheaf we see that
$$
\xymatrix{
q_0^*\mathcal{F}_T \ar[rr]_{q_{01}^*c} \ar[rd]_{q_{02}^*c} & &
q_1^*\mathcal{F}_T \ar[ld]^{q_{12}^*c} \\
& q_2^*\mathcal{F}_T
}
$$
is a commutative diagram of $\mathcal{O}_{T''}$-modules. For
$s \in \Gamma(T, \mathcal{F}_T)$ we have
$c(p_0^*s) = p_1^*s - \nabla(s)$. Write
$\nabla(s) = \sum p_1^*s_i \cdot \omega_i$ where $s_i$ is a local section
of $\mathcal{F}_T$ and $\omega_i$ is a local section of $\Omega_{T/S, \delta}$.
We think of $\omega_i$ as a local section of the structure
sheaf of $\mathcal{O}_{T'}$ and hence we write product instead of tensor
product. On the one hand
\begin{align*}
q_{12}^*c \circ q_{01}^*c(q_0^*s) & = 
q_{12}^*c(q_1^*s - \sum q_1^*s_i \cdot q_{01}^*\omega_i) \\
& =
q_2^*s - \sum q_2^*s_i \cdot q_{12}^*\omega_i -
\sum q_2^*s_i \cdot q_{01}^*\omega_i +
\sum q_{12}^*\nabla(s_i) \cdot q_{01}^*\omega_i
\end{align*}
and on the other hand
$$
q_{02}^*c(q_0^*s) = q_2^*s - \sum q_2^*s_i \cdot q_{02}^*\omega_i.
$$
From the formulae of Remark \ref{remark-second-order-thickening} we see
that
$q_{01}^*\omega_i + q_{12}^*\omega_i - q_{02}^*\omega_i = \text{d}\omega_i$.
Hence the difference of the two expressions above is
$$
\sum q_2^*s_i \cdot \text{d}\omega_i -
\sum q_{12}^*\nabla(s_i) \cdot q_{01}^*\omega_i
$$
Note that
$q_{12}^*\omega \cdot q_{01}^*\omega' = \omega' \wedge \omega =
- \omega \wedge \omega'$ by the definition of the multiplication on
$\mathcal{O}_{T''}$. Thus the expression above is $\nabla^2(s)$ viewed
as a section of the subsheaf $\mathcal{F}_T \otimes \Omega^2_{T/S, \delta}$ of
$q_2^*\mathcal{F}$. Hence we get the integrability condition.
\end{proof}
\end{document}
